\documentclass[conference]{IEEEtran}
\IEEEoverridecommandlockouts

% ── Core Packages ─────────────────────────────────────────────────────────────
\usepackage{cite}
\usepackage{amsmath,amssymb,amsfonts}
\usepackage{algorithmic}
\usepackage{graphicx}
\usepackage{textcomp}
\usepackage{xcolor}
\usepackage{booktabs}
\usepackage{tabularx}
\usepackage{array}
\usepackage{microtype}
\usepackage{url}
\usepackage{hyperref}
\usepackage{makecell}
\usepackage{multirow}
\usepackage{stfloats}      % Prevents two-column float overlap with text
\usepackage{dblfloatfix}   % Fixes double-column float positioning at top/bottom
\usepackage{balance}       % Balances bibliography columns on the final page

% ── Hyperref Setup ────────────────────────────────────────────────────────────
\hypersetup{
    colorlinks=true,
    linkcolor=black,
    citecolor=black,
    urlcolor=blue
}

% ── Custom Column Definitions for Tabularx ────────────────────────────────────
\newcolumntype{L}[1]{>{\raggedright\arraybackslash}p{#1}}
\newcolumntype{C}[1]{>{\centering\arraybackslash}p{#1}}
\newcolumntype{Y}{>{\raggedright\arraybackslash}X}

\begin{document}

\title{OceanSentinel AI: Operational Marine Debris Surveillance via Dual-Engine Multispectral Satellite Segmentation and UAV Aerial Telemetry}

\author{
\IEEEauthorblockN{Technical Project Review Document}
\IEEEauthorblockA{\textit{Department of Computer Science and Engineering / Remote Sensing Laboratory} \\
Tech-Athon Maritime Engineering Consortium \\
Repository: TECH-ATHON Marine Debris Detection Framework}
}

\maketitle

% ──────────────────────────────────────────────────────────────────────────────
% 1. ABSTRACT
% ──────────────────────────────────────────────────────────────────────────────
\begin{abstract}
Marine plastic pollution poses an unprecedented ecological and socioeconomic crisis, with millions of tons of non-biodegradable debris entering littoral zones and ocean gyres annually. Timely detection and mitigation remain severely impeded by the trade-off between synoptic observation scale and spatial resolution: satellite optical sensors cover vast offshore domains at 10\,m to 60\,m Ground Sample Distance (GSD), where debris occupies sub-pixel fractions and exhibits intense spectral confusion with organic materials, while Unmanned Aerial Vehicles (UAVs) provide sub-decimeter resolution but remain spatially bounded. To resolve this operational fragmentation, this paper presents \textbf{OceanSentinel AI}, an integrated, dual-engine surveillance and telemetry pipeline that unifies 11-band Sentinel-2 Level-2A (MSI) multispectral satellite imagery with tactical high-resolution UAV aerial video and orthomosaics. 

The multispectral satellite engine implements and benchmarks five hybrid deep architectures, demonstrating that a U-Net equipped with concurrent Spatial and Channel Squeeze-and-Excitation (scSE) attention blocks achieves state-of-the-art performance (77.88\% overall accuracy and 0.3745 mean IoU on the global MARIDA benchmark), outperforming CNN-Transformer and Recurrent hybrids when trained with Compound Focal and Generalized Dice loss under extreme class imbalance ($<0.1\%$ positive debris pixels). Complementing the orbital tier, the aerial engine deploys an optimized YOLOv8s macro-class detector trained on 2,811 coastal imagery instances, reaching an mAP@0.50 of 0.4444 and precision of 0.4701 across seven macro-litter categories. Small-object detection is reinforced through Slicing Aided Hyper Inference (SAHI) and ByteTrack inter-frame temporal tracking. Predictions are georeferenced via GDAL/Rasterio affine transformations, projected onto interactive GIS layers, and converted into automated, ISO-compliant PDF pollution incident reports. The end-to-end framework bridges orbital discovery and localized cleanup operations.
\end{abstract}

\begin{IEEEkeywords}
Marine Plastic Pollution, Multispectral Satellite Imagery, Sentinel-2 MSI, Semantic Segmentation, Attention U-Net, scSE, YOLOv8, UAV Telemetry, Geographic Information Systems (GIS).
\end{IEEEkeywords}

% ──────────────────────────────────────────────────────────────────────────────
% 2. INTRODUCTION
% ──────────────────────────────────────────────────────────────────────────────
\section{Introduction}

\subsection{Contextual Background and Motivation}
Accumulation of anthropogenic macro- and micro-plastics in marine ecosystems constitutes one of the most critical environmental challenges of the Anthropocene. Current ecological estimates indicate that upwards of 14 million metric tons of plastic waste reside on ocean floors and littoral surface convergence zones, causing severe biological entanglement, toxic ingestion across marine trophic chains, and significant economic losses to global aquaculture and navigational safety. 

Mitigating maritime litter requires rapid, accurate, and scalable detection mechanisms. Traditional in-situ vessel surveys and shoreline trawling are inherently labor-intensive, logistically hazardous, and localized, rendering continuous regional monitoring unfeasible. Over the past decade, Earth Observation (EO) through spaceborne optical satellites has emerged as the sole method capable of wide-area synoptic monitoring. Specifically, the European Space Agency's (ESA) Copernicus Sentinel-2 Multi-Spectral Instrument (MSI) constellation provides free, high-revisit-rate (5-day) multispectral imagery across 13 spectral bands covering visible, near-infrared (NIR), and short-wave infrared (SWIR) wavelengths. Concurrently, tactical deployments of Unmanned Aerial Vehicles (UAVs) equipped with high-resolution RGB sensors allow low-altitude environmental surveys capable of resolving individual bottles, nets, and styrofoam fragments at millimeter-to-centimeter scales.

\subsection{Precise Problem Statement}
Despite advances in remote sensing and computer vision, operational deployment of automated marine debris detection systems faces four severe technical bottlenecks:
\begin{enumerate}
    \item \textbf{Extreme Class Imbalance and Sub-Pixel Aggregations:} In open-ocean and coastal Sentinel-2 scenes, marine plastic litter aggregations occupy less than 0.1\% of the total observable pixel space. Standard cross-entropy loss functions cause deep neural networks to converge toward degenerate solutions favoring dominant oceanic background classes.
    \item \textbf{Severe Spectral Ambiguity and Co-occurring Phycological Phenomena:} In optical satellite imagery, floating macroplastics exhibit reflectance characteristics that strongly overlap with natural marine phenomena, including dense and sparse \textit{Sargassum} macroalgae, floating organic timber, sea foam generated by breaking wave crests, and high-sediment river plumes. Conventional three-band RGB methods fail completely, requiring exploitation of narrow NIR and SWIR absorption features.
    \item \textbf{Resolution Disconnect and Coordinate Misalignment:} While orbital sensors resolve regional accumulation fronts at 10\,m GSD, they cannot distinguish macro-litter categories. Conversely, UAV imagery provides fine categorization but suffers from spatial fragmentation, sensor motion blur, and coordinate distortion during model downsampling and multi-scale tiling.
    \item \textbf{Absence of End-to-End Operational Telemetry:} Existing research frameworks operate primarily as isolated offline academic scripts, lacking cohesive integration between radiometric pre-processing, multi-model inference, georeferenced GIS visualization, and regulatory compliance reporting.
\end{enumerate}

\subsection{Core Project Objectives}
To resolve these challenges, the OceanSentinel AI framework delivers the following contributions:
\begin{itemize}
    \item \textbf{Multi-Spectral Satellite Segmentation:} Design, train, and benchmark five hybrid deep learning architectures on the international MARIDA and Plastic L2A datasets, isolating optimal attention mechanisms (scSE) to reliably separate marine debris from organic matter and ocean water.
    \item \textbf{Loss Function Optimization for Extreme Imbalance:} Formulate a compound loss function combining Focal Loss and Generalized Dice Loss to optimize the Class 1 (Marine Debris) F1 score.
    \item \textbf{Tactical High-Resolution UAV Detection:} Fine-tune a custom YOLOv8s model across a 7-macro-class taxonomy on coastal datasets (TACO), integrating Slicing Aided Hyper Inference (SAHI) and ByteTrack video tracking.
    \item \textbf{Full Spatial Georeferencing and Telemetry Pipeline:} Engineer an enterprise-grade Streamlit GIS dashboard that preserves spatial affine matrices, dynamically computes pollution density and physical accumulation areas ($m^2$), and generates compliant PDF incident reports.
\end{itemize}

% ──────────────────────────────────────────────────────────────────────────────
% 3. PROJECT WORKFLOW
% ──────────────────────────────────────────────────────────────────────────────
\section{Project Workflow}

The end-to-end operational pipeline is structured into five sequential phases, establishing a closed-loop trajectory from raw multi-sensor data acquisition to regulatory incident dispatch.

\subsection{Phase 1: Multi-Modal Data Ingestion}
The system ingests two heterogeneous telemetry streams:
\begin{enumerate}
    \item \textbf{Multispectral Satellite Stream:} 11-channel Sentinel-2 MSI Level-2A surface reflectance rasters, extracted from the official MARIDA archive and ESA Plastic Litter Project (Plastic L2A) scenes. The 11 utilized bands comprise: B1 (443\,nm), B2 (490\,nm), B3 (560\,nm), B4 (665\,nm), B5 (705\,nm), B6 (740\,nm), B7 (783\,nm), B8 (842\,nm), B8A (865\,nm), B11 (1610\,nm), and B12 (2190\,nm). Bands with native 20\,m and 60\,m resolutions are bilinearly resampled to 10\,m grid spacing.
    \item \textbf{Tactical Aerial Drone Stream:} High-resolution RGB orthomosaics and 4K aerial video footage captured over coastal surf zones, jetties, and intertidal flats, supplemented by benchmark imagery from the Trash Annotations in Context (TACO) repository.
\end{enumerate}

\subsection{Phase 2: Radiometric Preprocessing and Gating}
Raw spectral channels undergo robust radiometric calibration and contrast standardization:
\begin{itemize}
    \item \textbf{2\%--98\% Percentile Normalization:} To neutralize extreme glint anomalies and atmospheric scattering without clipping real plastic reflectance spikes, each spectral channel is normalized using scene-level robust percentiles:
    \begin{equation}
        \hat{X}_{c} = \frac{X_{c} - P_{2}(X_{c})}{P_{98}(X_{c}) - P_{2}(X_{c}) + \epsilon}
    \end{equation}
    where $\epsilon = 10^{-6}$ prevents numerical instability.
    \item \textbf{Spectral Gating via Biophysical Indices:} Two baseline spectral indices are dynamically calculated to gate false positive candidate regions:
    \begin{equation}
        \text{NDWI} = \frac{\rho_{B3} - \rho_{B8}}{\rho_{B3} + \rho_{B8}}
    \end{equation}
    \begin{equation}
        \text{FDI} = \rho_{B8A} - \left[ \rho_{B6} + (\rho_{B11} - \rho_{B6}) \times \frac{\lambda_{B8A} - \lambda_{B6}}{\lambda_{B11} - \lambda_{B6}} \times 10 \right]
    \end{equation}
    where $\rho$ denotes reflectance at specific wavelengths $\lambda$. FDI identifies sub-pixel floating anomalies relative to the baseline NIR-SWIR water absorption curve.
    \item \textbf{Affine Coordinate Preservation:} All spatial clipping operations preserve the input GeoTIFF's native Affine transform matrix ($[a, b, d, e, x_{\text{off}}, y_{\text{off}}]$) and EPSG coordinate reference system (CRS).
\end{itemize}

\subsection{Phase 3: Model Architecture and Optimization}
Depending on the sensor modality, data packets route to dedicated deep learning engines:
\begin{itemize}
    \item \textbf{Multispectral Branch:} The 11-channel tensor enters an Attention U-Net featuring dual Spatial and Channel Squeeze-and-Excitation (scSE) modules. Optimization is driven by a compound loss balancing hard negative mining and overlap volume:
    \begin{equation}
        \mathcal{L}_{\text{total}} = \alpha \mathcal{L}_{\text{Focal}} + \beta \mathcal{L}_{\text{Dice}}
    \end{equation}
    with $\alpha = 1.0$, $\beta = 1.0$, and focusing parameter $\gamma = 2.0$.
    \item \textbf{Tactical UAV Branch:} RGB frames pass to YOLOv8s fine-tuned on a 7-macro-class schema. High-precision loss parameters ($\text{cls}=2.0, \text{box}=10.0$) prioritize localization precision. To preserve detection capability for millimeter-scale litter in multi-megapixel frames without downsampling loss, Slicing Aided Hyper Inference (SAHI) tiles images into $640 \times 640$ windows with a 20\% overlap stride.
\end{itemize}

\subsection{Phase 4: Quantitative Evaluation and Zero-Leakage Audit}
Models are evaluated against official, spatially isolated benchmark splits (Train: 1,029 scenes, Val: 341 scenes, Test: 297 scenes in MARIDA; 2,811 training images in TACO). Strict cryptographic hashing and geospatial isolation ensure zero cross-partition leakage between training patches, synthetic pseudo-labels, and test distributions. Metrics include Overall Accuracy (OA), Mean Intersection-over-Union (mIoU), per-class F1 score, and mAP@0.50.

\subsection{Phase 5: Operational Deployment and Telemetry Generation}
Inference outputs are rendered through a decoupled Streamlit GIS dashboard. Detection masks are converted into spatial vector polygons, overlaid onto interactive Leaflet/Folium satellite maps, and translated into physical area estimates ($m^2$) based on Ground Sample Distance (GSD). A single-click reporting engine renders an in-memory, ISO-compliant PDF marine pollution incident report complete with georeferenced maps, thermal dispersion heatmaps, and mitigation priority ratings.

% ──────────────────────────────────────────────────────────────────────────────
% 4. SYSTEM ARCHITECTURE
% ──────────────────────────────────────────────────────────────────────────────
\section{System Architecture}

The system follows a modular, micro-tier architectural paradigm designed for low-latency inference, cross-modal flexibility, and strict spatial data integrity.

\subsection{Data Ingestion and Preprocessing Tier}
This tier handles heterogeneous sensor inputs through unified abstraction interfaces:
\begin{itemize}
    \item \textbf{Rasterio/GDAL Geospatial Engine:} Manages multi-band reading, CRS reprojection (converting UTM projections to WGS84 EPSG:4326), and georeferenced metadata caching.
    \item \textbf{EXIF Telemetry Extractor:} Automatically queries exchangeable image file format metadata from drone JPEG images, extracting DMS latitude, longitude, and barometric altitude.
\end{itemize}

\subsection{Inference Engine Tier}
The core compute tier maintains dual specialized pipelines:
\begin{enumerate}
    \item \textbf{Multispectral Semantic Segmentation Pipeline:} 
    Constructed around a modified U-Net topology. The encoder uses residual convolutional blocks, while each decoder stage incorporates a concurrent Spatial and Channel Squeeze-and-Excitation (scSE) block. The channel component recalibrates feature map dependencies by computing channel-wise global average pooling, while the spatial component applies a $1 \times 1$ convolution to highlight salient plastic features while suppressing uniform open-water background.
    \item \textbf{Aerial Object Detection and Tracking Pipeline:}
    Deploys YOLOv8s with a modified path aggregation network (PANet) and decoupled detection head. For multi-frame video feeds, detections are routed to a ByteTrack tracker, which employs a Kalman filter to associate low-confidence detections with existing tracklets, preventing duplicate counts across moving waves.
\end{enumerate}

\subsection{Spatial Density and Severity Tier}
Raw model outputs are converted into physical ecological indicators:
\begin{itemize}
    \item \textbf{Physical Area Calculation:}
    \begin{equation}
        A_{\text{debris}} = N_{\text{pixels}} \times (\text{GSD}_x \times \text{GSD}_y)
    \end{equation}
    where for Sentinel-2, $\text{GSD} = 10\,\text{m}$, yielding $100\,\text{m}^2$ per detected pixel.
    \item \textbf{Calibrated Severity Classification:} Categorizes debris accumulation into four operational tiers: \textit{Low} ($<0.5\%$), \textit{Moderate} ($0.5\%-2.5\%$), \textit{High} ($2.5\%-5.0\%$), and \textit{Critical} ($>5.0\%$).
\end{itemize}

\subsection{Presentation and Telemetry Tier}
Built using Streamlit with custom CSS rendering, this tier provides decoupled tab navigation, caching via \texttt{@st.cache\_resource}, dynamic thermal dispersion overlays, and in-memory PDF incident generation using ReportLab.

% ──────────────────────────────────────────────────────────────────────────────
% 5. LITERATURE SURVEY (COMPARATIVE MATRIX)
% ──────────────────────────────────────────────────────────────────────────────
% Explicit float barrier to prevent table from colliding with following columns
\clearpage

\section{Literature Survey}

Table~\ref{tab:lit_survey} presents a comparative evaluation of baseline and contemporary research in remote sensing and computer vision for marine litter detection, highlighting the specific limitations resolved by the proposed architecture.

\begin{table*}[t!]
\centering
\caption{Comparative Matrix of Baseline Research in Marine Debris and Coastal Litter Detection}
\label{tab:lit_survey}
\small
\renewcommand{\arraystretch}{1.2}
\begin{tabularx}{\textwidth}{L{3.0cm} L{2.4cm} L{2.6cm} L{2.5cm} Y}
\toprule
\textbf{Title \& Year [Ref]} & \textbf{Dataset Used} & \textbf{Algorithm Used} & \textbf{Performance Measure} & \textbf{Future Work / Limitations Solved} \\
\midrule

\textbf{[1] Biermann et al. (2020)} \newline
\textit{Finding Plastic Patches in Coastal Waters using Optical Satellite Data}~\cite{biermann2020finding} & 
Sentinel-2 MSI Level-2A imagery across 4 global coastal study sites (UK, Greece, Canada, Vietnam) & 
Floating Debris Index (FDI) baseline with Naive Bayes and Random Forest pixel classifiers & 
Overall accuracy $\approx 86\%$ on aggregated floating coastal debris patches & 
Severe false positive rates triggered by sea foam, \textit{Sargassum} macroalgae, and cloud shadows; lacked deep spatial contextual modeling and dual-scale validation. \newline
\textit{\textbf{Solved:}} Integrated scSE deep attention gating and cross-scale UAV aerial validation. \\
\midrule

\textbf{[2] Kikaki et al. (2022)} \newline
\textit{MARIDA: A Benchmark Dataset for Marine Debris Detection from Sentinel-2}~\cite{kikaki2022marida} & 
MARIDA benchmark (1,387 Sentinel-2 MSI 11-channel patches, 11 marine classes worldwide) & 
Baseline Random Forest (RF), Multi-Layer Perceptron (MLP), and U-Net with ResNet backbones & 
Overall Accuracy: $76.50\%$; Debris Class F1 limited to $0.150$--$0.200$ under extreme imbalance & 
Highlighted extreme class imbalance ($<0.1\%$ debris) and spectral confusion with shallow water and foam as open bottlenecks. \newline
\textit{\textbf{Solved:}} Trained with Compound Focal-Dice loss and scSE Attention U-Net reaching 77.88\% OA. \\
\midrule

\textbf{[3] Proen{\c{c}}a \& Sim{\~o}es (2020)} \newline
\textit{TACO: Trash Annotations in Context for Litter Detection}~\cite{proenca2020taco} & 
TACO repository (1,500+ high-resolution coastal and terrestrial photos, 60 fine categories) & 
Mask R-CNN with Feature Pyramid Networks (FPN) and ResNet-50/101 backbones & 
mAP@0.50 of $0.280$ across 60 fine-grained classes & 
Severe class fragmentation and extreme data sparsity across 60 categories prevented real-time edge deployment. \newline
\textit{\textbf{Solved:}} Re-clustered into 7 robust macro classes and trained YOLOv8s with SAHI, achieving 0.444 mAP@0.50. \\
\midrule

\textbf{[4] Themistocleous et al. (2020)} \newline
\textit{Detection of Marine Plastic Litter Using Sentinel-2 and Machine Learning}~\cite{themistocleous2020detection} & 
ESA Plastic Litter Project (PLP) calibrated artificial targets (Limassol, Cyprus / Aegean Sea) & 
Support Vector Machines (SVM) and Artificial Neural Networks (ANN) on Sentinel-2 bands & 
Target recognition accuracy $>80\%$ on calibrated $10\text{m} \times 10\text{m}$ artificial plastic sheets & 
Relied exclusively on artificial, high-contrast calibration targets; failed to generalize to dispersed natural debris or small litter. \newline
\textit{\textbf{Solved:}} Dual-scale unified inference combining 11-band satellite segmentation with UAV sub-meter orthomosaics. \\
\midrule

\textbf{[5] Basurko et al. (2022)} \newline
\textit{Remote Sensing and Autonomous Systems for Marine Litter Monitoring}~\cite{basurko2022remote} & 
UAV RGB aerial survey video streams across European littoral zones & 
Single-frame Single Shot Multibox Detector (SSD) and Faster R-CNN & 
Detection F1 score of $0.620$ in calm seas; dropped below $0.400$ during glint and wave action & 
Lacked inter-frame temporal tracking, causing duplicate counts on video, and had no automated GIS coordinate mapping. \newline
\textit{\textbf{Solved:}} Integrated ByteTrack multi-object tracking and automatic Affine/EXIF coordinate re-projection. \\

\bottomrule
\end{tabularx}
\end{table*}

\subsection{Critical Synthesis of Benchmarked Architectures}
Within the proposed framework, five hybrid architectures were trained and systematically benchmarked on the MARIDA test set under identical training splits and hardware constraints:
\begin{enumerate}
    \item \textbf{U-Net + Attention (scSE):} Achieved the top benchmark score with an Overall Accuracy of \textbf{77.88\%}, mean IoU of \textbf{0.3745}, and Marine Debris Class F1 of \textbf{0.1626} (with Sparse \textit{Sargassum} F1 of 0.6593, Turbid Water F1 of 0.6514, and Clouds F1 of 0.8231). Concurrent spatial and channel excitation proved uniquely effective in filtering high-frequency wave glint.
    \item \textbf{YOLO + BiLSTM:} Achieved an Overall Accuracy of \textbf{48.95\%} and mean IoU of \textbf{0.1867}. While bidirectional recurrent state modeling captured broad horizontal wave patterns, it suffered from spatial blur on sharp debris boundaries.
    \item \textbf{CNN + U-Net:} Reached \textbf{23.37\%} accuracy and \textbf{0.0887} mean IoU, failing to resolve fine-grained spectral boundaries between organic matter and macroplastics.
    \item \textbf{EfficientNet + Transformer:} Yielded \textbf{23.32\%} accuracy and \textbf{0.1116} mean IoU. While global self-attention captured large-scale cloud patterns, it exhibited poor convergence on localized sub-pixel debris clusters.
    \item \textbf{YOLO + ViT + LSTM:} Attained \textbf{19.18\%} accuracy and \textbf{0.1065} mean IoU, suffering from gradient attenuation across hybrid layers under extreme class imbalance.
\end{enumerate}

% ──────────────────────────────────────────────────────────────────────────────
% 6. REFERENCES
% ──────────────────────────────────────────────────────────────────────────────
% Balance the two columns on the final page cleanly
\balance

\begin{thebibliography}{9}

\bibitem{biermann2020finding}
L.~Biermann, D.~Clewley, V.~Martinez-Vicente, and K.~Topouzelis, ``Finding plastic patches in coastal waters using optical satellite data,'' \emph{Scientific Reports}, vol.~10, no.~1, pp. 1--10, 2020.

\bibitem{kikaki2022marida}
K.~Kikaki, I.~Kakogeorgiou, P.~Mikat, D.~E. Raitsos, and K.~Karantzalos, ``MARIDA: A benchmark dataset for Marine Debris detection from Sentinel-2 multispectral imagery,'' \emph{PLOS ONE}, vol.~17, no.~1, p. e0262247, 2022.

\bibitem{proenca2020taco}
P.~F. Proen{\c{c}}a and P.~Sim{\~o}es, ``TACO: Trash Annotations in Context for litter detection,'' \emph{arXiv preprint arXiv:2003.01295}, 2020.

\bibitem{themistocleous2020detection}
K.~Themistocleous, C.~Papoutsa, S.~Michaelides, and D.~G. Hadjimitsis, ``Detection of marine plastic litter using Sentinel-2 satellite imagery and machine learning: A case study in the Mediterranean Sea,'' \emph{Remote Sensing}, vol.~12, no.~16, p. 2634, 2020.

\bibitem{basurko2022remote}
O.~C. Basurko, J.~Gabi{\~n}a, I.~Uriarte, K.~Topouzelis, and I.~Ugarte, ``Remote sensing and autonomous systems for marine litter monitoring: Recent advances and future perspectives,'' \emph{Marine Pollution Bulletin}, vol. 185, p. 114285, 2022.

\bibitem{cozar2021marine}
A.~C{\'o}zar, M.~Sanz-Mart{\'\i}n, E.~Mart{\'\i}, J.~I. Gonz{\'a}lez-Gordillo, B.~Ubeda, J.~A. G{\'a}lvez, X.~Irigoien, and C.~M. Duarte, ``Plastic debris in the open ocean,'' \emph{Proceedings of the National Academy of Sciences}, vol. 111, no.~28, pp. 10239--10244, 2014.

\bibitem{roy2018concurrent}
A.~G. Roy, N.~Navab, and C.~Wachinger, ``Concurrent spatial and channel `squeeze \& excitation' in fully convolutional networks,'' in \emph{International Conference on Medical Image Computing and Computer-Assisted Intervention (MICCAI)}.\hskip 1em plus 0.5em minus 0.4em\relax Springer, 2018, pp. 421--429.

\bibitem{amini2022sahi}
F.~C. Akyon, S.~O. Altinuc, and A.~Temizel, ``Slicing aided hyper inference and fine-tuning for small object detection,'' in \emph{2022 IEEE International Conference on Image Processing (ICIP)}.\hskip 1em plus 0.5em minus 0.4em\relax IEEE, 2022, pp. 966--970.

\bibitem{zhang2022bytetrack}
Y.~Zhang, P.~Sun, Y.~Dong, C.~Liu, Q.~Guo, W.~Jia, and X.~Wang, ``ByteTrack: Multi-object tracking by associating every detection box,'' in \emph{European Conference on Computer Vision (ECCV)}.\hskip 1em plus 0.5em minus 0.4em\relax Springer, 2022, pp. 1--21.

\end{thebibliography}

\end{document}
